\documentclass[11pt,a4paper]{article}
\usepackage[margin=27mm]{geometry}
\usepackage[T1]{fontenc}
\usepackage{lmodern,amsmath,amssymb,amsthm,mathtools,microtype}
\usepackage[colorlinks=true,linkcolor=blue,citecolor=blue,urlcolor=blue]{hyperref}
\newtheorem{theorem}{Theorem}
\newtheorem{lemma}[theorem]{Lemma}
\newtheorem{corollary}[theorem]{Corollary}
\newtheorem{proposition}[theorem]{Proposition}
\newcommand{\Tr}{\operatorname{Tr}}
\newcommand{\rank}{\operatorname{rank}}
\newcommand{\Herm}{\operatorname{Herm}}
\newcommand{\id}{\operatorname{id}}
\newcommand{\F}{\mathrm F}
\newcommand{\C}{\mathbb C}
\newcommand{\cS}{\mathcal S}
\newcommand{\cK}{\mathcal K}
\newcommand{\cU}{\mathcal U}
\newcommand{\norm}[1]{\left\lVert#1\right\rVert_{\F}}
\title{Sharp marginal-erasure bounds and positive completion\\of nonunital quantum channels}
\author{Lluis Eriksson}
\date{24 October 2026}
\begin{document}
\maketitle
\begin{abstract}
We determine the exact Frobenius conditioning of deleting the traceless
part of one local matrix algebra on $\C^m\otimes\C^k$.
For $1\le r<m$, the inverse constants are
$\sqrt{m/(m-r)}$ for two states of rank at most $r$, and
$\sqrt{(m+r)/(m-r)}$ for one such state against an arbitrary density
matrix. Both are attained, independently of $k\ge2$; at rank $m$ exact
collisions occur. The one-sided estimate follows from a scalar
optimization combining positivity, trace control and the partial-trace
norm. It yields a noise and spectral-tail bound for convex positive
completion. For normalized Choi matrices, the retained coordinates are
the action of a trace-preserving channel on traceless inputs.
Consequently this action identifies every channel of Kraus rank below
its input dimension among all channels, without a unitality assumption.
For square channels we express the same conclusion through a complete
exact unitary $2$-design of average gate fidelities. The rank threshold
is sharp, whereas optimal conditioning within the trace-preserving
subclass is left open. The fidelity frame identities and general
kernel/inertia method are prior work; no compressed-sampling or new
randomized-benchmarking noise guarantee is asserted.
\end{abstract}

\section{Question, contribution and antecedents}
Average gate fidelities with unitary reference gates do not linearly
determine an arbitrary quantum channel: they miss its translation, or
action on the maximally mixed input. Positivity can nevertheless make
that missing part unique when the channel has sufficiently small Kraus
rank. We quantify this mechanism by first solving an unrestricted
state-space problem exactly.

Let $\cS_{m,k}$ denote the density matrices on
$\C^m\otimes\C^k$, and set
\[
 \cK=\{I_m\otimes B:B\in\Herm(\C^k),\ \Tr B=0\}.
\]
Here $\cK$ is the traceless Hermitian part of a local algebra, not itself
an algebra under multiplication. Let $P$ be its Frobenius-orthogonal
projection and $Q=\id-P$. Our principal result computes the largest
ratio $\norm{\rho-\sigma}/\norm{Q(\rho-\sigma)}$ under either one or
two rank constraints. This distinction matters: positivity-constrained
reconstruction usually produces a competitor with unrestricted rank.

The reduction of low-rank identification to the inertia of a measurement
kernel is classical in state tomography
\cite{carmeli} and channel tomography \cite{gutoski}.
In particular, Gutoski and Johnston construct small families of
interactive observables identifying low-Kraus-rank channels among all
channels. We do not improve their observable counts; we study the
conditioning of a particular fixed omitted subspace.
Kimmel et al.\ recover the unital block using randomized benchmarking
\cite{kimmel}; their single-qubit complete-positivity lemma has a
corrected proof \cite{erratum}. Roth et al.\ provide a unitary-design
frame for unital trace-preserving maps and compressed recovery
guarantees for quantum gates \cite{roth}. We use the classical
second-moment identity explicitly and separately from our completion
inequality.

The proposed contribution is the pair of exact local-erasure constants,
their state attainers, and the resulting quantitative completion
statement without a unitality promise. The elementary proof is included
in full. These comparisons delimit the claim; they do not certify
worldwide priority.

\section{Exact conditioning in state space}
All matrix spaces below are real Hilbert spaces of Hermitian matrices
with inner product $\Tr(XY)$. The Frobenius norm is $\norm X$.
A zero denominator and nonzero numerator mean an infinite ratio.
Write $\cS^{(r)}_{m,k}$ for states of rank at most $r$.

\begin{theorem}[Exact local-erasure constants]\label{thm:sharp}
Let $m\ge1$, $k\ge2$ and $r\ge1$ be integers. For $1\le r<m$,
\begin{align}
 \sup_{\substack{\rho,\sigma\in\cS^{(r)}_{m,k}\\\rho\ne\sigma}}
 \frac{\norm{\rho-\sigma}}{\norm{Q(\rho-\sigma)}}
 &=\sqrt{\frac{m}{m-r}}, \label{eq:balanced}\\
 \sup_{\substack{\rho\in\cS^{(r)}_{m,k},\ \sigma\in\cS_{m,k}\\\rho\ne\sigma}}
 \frac{\norm{\rho-\sigma}}{\norm{Q(\rho-\sigma)}}
 &=C_{m,r}:=\sqrt{\frac{m+r}{m-r}}. \label{eq:one}
\end{align}
Both suprema are attained. For $r\ge m$ both are infinite.
Thus $m=1$ has no finite positive-rank regime.
\end{theorem}

We first record the projection and the standard inertia reduction.
The notation $n_+(D)$ and $n_-(D)$ counts positive and negative
eigenvalues with multiplicity.

\begin{lemma}\label{lem:projection}
For any Hermitian $D$,
\[
 PD=\frac{I_m}{m}\otimes
 \left(\Tr_A D-\frac{\Tr D}{k}I_k\right).
\]
In particular, for $\Tr D=0$,
\begin{equation}\label{eq:projection}
 PD=\frac{I_m}{m}\otimes\Tr_A D,\qquad
 \norm{PD}^2=\frac1m\norm{\Tr_A D}^2.
\end{equation}
For arbitrary matrices $X$, $\norm{\Tr_A X}\le\sqrt m\,\norm X$.
If $D=\rho-\sigma$ with $\rank\rho\le r$, then $n_+(D)\le r$.
If also $\rank\sigma\le r$, then $n_-(D)\le r$.
\end{lemma}
\begin{proof}
The proposed $PD$ belongs to $\cK$, and its difference from $D$ is
orthogonal to every $I_m\otimes B$ with $\Tr B=0$, by the defining
identity of partial trace. This proves the projection formula.
For the norm estimate, in a product basis Cauchy--Schwarz gives
\[
 \sum_{i,j}\left|\sum_a X_{ai,aj}\right|^2
 \le m\sum_{a,i,j}|X_{ai,aj}|^2\le m\norm X^2.
\]
If a positive eigenspace of $\rho-\sigma$ had dimension greater than
$r$, it would meet $\ker\rho$ nontrivially. On that kernel its quadratic
form is nonpositive, a contradiction. Changing sign proves the last
claim.
\end{proof}

\begin{proof}[Proof of Theorem~\ref{thm:sharp}]
Let $D\ne0$ be Hermitian with $\Tr D=0$, and write its spectral
decomposition as $D=D_+-D_-$ with orthogonal positive parts. Set
\[
 \tau=\Tr D_+=\Tr D_->0,\qquad
 a^2=\norm{D_+}^2,\quad b^2=\norm{D_-}^2.
\]
If $n_+(D)\le r$, then $a^2\ge\tau^2/r$. The positive matrices
$R=\Tr_A D_+$ and $S=\Tr_A D_-$ both have trace $\tau$.
Positivity implies $\Tr(RS)\ge0$ and
$\norm R^2,\norm S^2\le\tau^2$. Lemma~\ref{lem:projection} also gives
$\norm S^2\le m b^2$. Consequently
\begin{equation}\label{eq:key}
 \norm{\Tr_A D}^2
 =\norm{R-S}^2
 \le \tau^2+\min\{\tau^2,m b^2\}.
\end{equation}
No commutation or orthogonality of the two marginals is assumed.

Put $t=b^2/\tau^2$. Equations~\eqref{eq:projection}--\eqref{eq:key} yield
\begin{equation}\label{eq:scalar}
 \frac{\norm{PD}^2}{\norm D^2}
 \le \frac{1/m+\min\{1/m,t\}}{1/r+t}
 \le \frac{2r}{m+r}.
\end{equation}
Indeed the middle expression is increasing for $0\le t\le1/m$ and
decreasing for $t\ge1/m$, since $r<m$. Its maximum is at $t=1/m$.
Orthogonality of $P$ and $Q$ now proves \eqref{eq:one} as an upper
bound for every difference in its stated domain.

For sharpness, choose a rank-$r$ input projector $R_0$ and orthogonal
output unit vectors $u,v$. The states
\begin{equation}\label{eq:one-attain}
 \rho=\frac{R_0}{r}\otimes|u\rangle\langle u|,
 \qquad
 \sigma=\frac{I_m}{m}\otimes|v\rangle\langle v|
\end{equation}
have ranks $r,m$. For their difference,
\[
 \norm D^2=\frac1r+\frac1m,\quad
 \norm{PD}^2=\frac2m,\quad
 \norm{QD}^2=\frac1r-\frac1m.
\]
They attain \eqref{eq:one} and all bounds in \eqref{eq:scalar}.

If both inertia indices are at most $r$, then
$a^2,b^2\ge\tau^2/r$. Instead of \eqref{eq:key} use
\[
 \norm{\Tr_A D}^2\le2\tau^2\le r(a^2+b^2).
\]
The hidden fraction is at most $r/m$, which proves the upper bound in
\eqref{eq:balanced}. It is attained by
\[
 \rho=\frac{R_0}{r}\otimes|u\rangle\langle u|,\qquad
 \sigma=\frac{R_0}{r}\otimes|v\rangle\langle v|,
\]
for which $\norm D^2=2/r$ and $\norm{PD}^2=2/m$.
Finally, take $R_0=I_m$ and $r=m$. These distinct rank-$m$ states
have $QD=0$. They remain admissible for every $r\ge m$, establishing
the infinite regime, including $m=1$.
\end{proof}

Every nonzero member of $\cK$ has at least $m$ positive and $m$
negative eigenvalues. Thus the qualitative threshold also follows
immediately from the classical kernel/inertia criterion. The estimates
above quantify the angle between that kernel and the relevant
difference cones, rather than merely showing they do not intersect.

\section{Approximate rank and convex completion}
For a density matrix $\rho$ with eigenvalues in decreasing order, define
$\tau_r(\rho)=\sum_{i>r}\lambda_i(\rho)$.

\begin{proposition}[Spectral tails]\label{prop:tail}
For $1\le r<m$, every pair of density matrices $\rho,\sigma$ obeys
\begin{equation}\label{eq:tail}
 \norm{\sigma-\rho}
 \le C_{m,r}\norm{Q(\sigma-\rho)}
       +\sqrt2(C_{m,r}+1)\tau_r(\rho).
\end{equation}
Let $\mathcal F$ be any compact convex subset of density matrices
containing $\rho$. If $y=Q\rho+e$, $\norm e\le\epsilon$, then every
minimizer $\widehat\rho$ of $\norm{Q\sigma-y}$ over $\sigma\in\mathcal F$
satisfies
\begin{equation}\label{eq:recovery}
 \norm{\widehat\rho-\rho}
 \le2C_{m,r}\epsilon+\sqrt2(C_{m,r}+1)\tau_r(\rho).
\end{equation}
\end{proposition}
\begin{proof}
Let $s=1-\tau_r(\rho)>0$, and let $\rho_r$ be the top-$r$ spectral
truncation of $\rho$ divided by $s$. At an eigenvalue tie any spectral
choice is allowed. Then
\[
 \norm{\rho-\rho_r}^2
 =\tau_r(\rho)^2\sum_{i\le r}\left(\frac{\lambda_i}{s}\right)^2
   +\sum_{i>r}\lambda_i^2
 \le2\tau_r(\rho)^2.
\]
Apply \eqref{eq:one} to $\rho_r,\sigma$, then use the triangle
inequality and the contraction property of $Q$ to obtain
\eqref{eq:tail}. A minimizer exists by compactness. Its residual is
at most that of the feasible truth, hence
$\norm{Q(\widehat\rho-\rho)}\le2\epsilon$.
Substitution proves \eqref{eq:recovery}.
\end{proof}

The truncation $\rho_r$ need not belong to $\mathcal F$. The proof is
valid because \eqref{eq:one} compares it with every density matrix.
This detail matters for Choi matrices: normalized spectral truncation
does not in general preserve trace preservation. For example, the
$3\to2$ channel with Kraus operators
\[
 K_0=\begin{pmatrix}1&0&0\\0&1&0\end{pmatrix},\qquad
 K_1=\begin{pmatrix}0&0&1\\0&0&0\end{pmatrix}
\]
has Choi eigenvalues $2/3,1/3$. Its normalized top-one truncation is
the state of $(|00\rangle+|11\rangle)/\sqrt2$, whose input marginal is
$\operatorname{diag}(1/2,1/2,0)$ rather than $I_3/3$.
Convexity makes the optimization a convex problem; the proof of the
error bound uses compactness and feasibility, without imposing a rank
constraint on the minimizer.

\section{Channel completion from the traceless-input action}
Let $\Phi:M_m\to M_k$ be completely positive and trace preserving
(CPTP). Use input--output tensor order and the normalized Choi matrix
\[
 J_\Phi=(\id\otimes\Phi)(|\Omega_m\rangle\langle\Omega_m|),\qquad
 |\Omega_m\rangle=\frac1{\sqrt m}\sum_{a=0}^{m-1}|a,a\rangle.
\]
Then $J_\Phi\ge0$, $\Tr J_\Phi=1$, $\Tr_{\rm out}J_\Phi=I_m/m$,
and $\rank J_\Phi$ is the minimum Kraus rank.

Choose Frobenius-orthonormal bases $F_i$ and $G_j$ of the real
spaces of traceless Hermitian input and output matrices. Define
\[
 T(\Phi)_{ji}=\Tr[G_j\Phi(F_i)].
\]
These bases can contain imaginary entries. The input transpose in
\begin{equation}\label{eq:choi-coordinate}
 \Tr[J_\Phi(F_i^{\mathsf T}\otimes G_j)]
 =\frac1m T(\Phi)_{ji}
\end{equation}
is essential. It follows by expanding
$J_\Phi=m^{-1}\sum_{a,b}|a\rangle\langle b|\otimes
\Phi(|a\rangle\langle b|)$.

For two TP maps the identity and input-only coefficients of
$D=J_\Phi-J_\Psi$ vanish. The projection $P$ removes precisely
the output-only traceless coefficients, so orthonormal expansion gives
\begin{equation}\label{eq:choi-norm}
 \norm{Q(J_\Phi-J_\Psi)}
 =\frac1m\norm{T(\Phi)-T(\Psi)}.
\end{equation}
Thus $T$ is exactly the retained nonconstant Choi information. The
omitted translation is $\Phi(I_m/m)$.

\begin{theorem}[Low-Kraus-rank completion]\label{thm:channel}
If $\rank J_\Phi\le r<m$, then for every CPTP $\Psi:M_m\to M_k$,
\begin{equation}\label{eq:channel}
 \norm{J_\Phi-J_\Psi}
 \le\frac{C_{m,r}}m\norm{T(\Phi)-T(\Psi)}.
\end{equation}
In particular $T(\Phi)$ identifies $\Phi$ among all CPTP maps.
The uniform rank threshold $r<m$ is sharp for every $k\ge2$.
For $T_{\rm obs}=T(\Phi)+N$ with $\norm N\le\delta$, every least-squares
fit over all CPTP maps satisfies, for every $1\le r<m$,
\begin{equation}\label{eq:channel-noise}
 \norm{J_{\widehat\Phi}-J_\Phi}
 \le\frac{2C_{m,r}}m\delta+
       \sqrt2(C_{m,r}+1)\tau_r(J_\Phi).
\end{equation}
\end{theorem}
\begin{proof}
Combine Theorem~\ref{thm:sharp} and \eqref{eq:choi-norm}.
For the endpoint consider the pure-output replacers
$\Phi_u(X)=\Tr(X)|u\rangle\langle u|$.
Each has $T=0$ and Choi rank $m$; distinct output states give distinct
channels with identical data.

For the noisy result, the feasible Choi matrices form a compact convex
set with marginal $I_m/m$. Insert $T_{\rm obs}/m$ into the correlation
coordinates, together with the known identity coefficient and
zero input-only coefficients. This is $QJ_\Phi+e$ with
$\norm e\le\delta/m$.
The residual problem differs from the displayed least-squares problem
only by the factor $1/m$, so Proposition~\ref{prop:tail} applies.
\end{proof}

The state-space constants are exact; \eqref{eq:channel} is a valid upper
bound on the smaller fixed-marginal class, not an optimal channel
constant proved here. In \eqref{eq:one-attain} the first input marginal
is $R_0/r$, so that attainer is not TP for $r<m$.
For rectangular maps some rank classes are empty: $r$ Kraus operators
of size $k$ by $m$ require $rk\ge m$. For square maps Kraus rank one
means a unitary channel, whose translation is zero.

\subsection{Why one must not impose unitality}
For square dimensions let the conventional unital part be the
Hermiticity-preserving TP map
\begin{equation}\label{eq:unital}
 \Phi^{\rm u}(X)=\Phi(X)
       -\frac{\Tr X}{m}\bigl(\Phi(I_m)-I_m\bigr).
\end{equation}
It fixes $I_m$ and agrees with $\Phi$ on traceless inputs.

\begin{corollary}\label{cor:unital}
If $\Phi:M_m\to M_m$ is nonunital and has Kraus rank less than $m$,
then $\Phi^{\rm u}$ is not completely positive.
\end{corollary}
\begin{proof}
Otherwise $\Phi^{\rm u}$ would be a distinct CPTP channel with
$T(\Phi^{\rm u})=T(\Phi)$, contradicting Theorem~\ref{thm:channel}.
\end{proof}

For a concrete family, let $m\ge3$ and
\[
 K_0=\operatorname{diag}(1,\ldots,1,3/5),\qquad
 K_1=(4/5)|0\rangle\langle m-1|.
\]
The two linearly independent Kraus operators satisfy
$\sum_aK_a^\dagger K_a=I_m$ and define a nonunital channel of rank two.
Writing $B=\Phi(I_m)-I_m=(16/25)(|0\rangle\langle0|
-|m-1\rangle\langle m-1|)$, its unital-part Choi matrix is
$J_{\Phi^{\rm u}}=J_\Phi-I_m\otimes B/m^2$.
For any $1\le i\le m-2$,
\[
 \langle i,0|J_{\Phi^{\rm u}}|i,0\rangle=-\frac{16}{25m^2}<0.
\]
This illustrates the failure of an unjustified unital constraint.
It is not a claim that higher-dimensional failure of complete
positivity was previously unknown. There is no conflict with the
corrected qubit lemma of \cite{erratum}: below the qubit threshold
the only square channels are unitary, and the interior index in this
example does not exist for $m=2$.

\section{Complete average-gate-fidelity data}
In this section input and output dimensions both equal $m\ge2$.
Let $U_1,\ldots,U_N$ be a finite exact unitary $2$-design with uniform
weights, and put
\[
 f_\ell(\Phi)=F_{\rm avg}(U_\ell,\Phi)
 =\int\langle\psi|U_\ell^\dagger
       \Phi(|\psi\rangle\langle\psi|)U_\ell|\psi\rangle\,d\psi,
 \qquad
 \|z\|_{\rm rms}^2=\frac1N\sum_\ell|z_\ell|^2 .
\]
For a weighted design the same statements use positive weights
$w_\ell$ with $\sum_\ell w_\ell=1$ and the matched norm
$\|z\|_{{\rm rms},w}^2=\sum_\ell w_\ell|z_\ell|^2$.
No random subset of this design is used.

\begin{lemma}[Classical design frame, with normalization]\label{lem:agf}
For any two CPTP maps on $M_m$,
\begin{equation}\label{eq:frame}
 \norm{Q(J_\Phi-J_\Psi)}
 =(m+1)\sqrt{m^2-1}\,\|f(\Phi)-f(\Psi)\|_{\rm rms}.
\end{equation}
\end{lemma}
\begin{proof}
Set $p=m^2-1$ and use the same orthonormal traceless basis $F_i$ on
input and output. Let $O(U)_{ji}=\Tr[F_jUF_iU^\dagger]$.
The standard pure-state second moment
$\int|\psi\rangle\langle\psi|^{\otimes2}d\psi
=(I+\mathrm{SWAP})/[m(m+1)]$ gives, for a TP map,
\begin{equation}\label{eq:favg}
 F_{\rm avg}(U,\Phi)=\frac1m+
 \frac{\Tr[O(U)^{\mathsf T}T(\Phi)]}{m(m+1)}.
\end{equation}
For completeness the design second moment follows from
\[
 \int(UF_iU^\dagger)\otimes(UF_bU^\dagger)\,dU
 =\frac{\delta_{ib}}{p}\left(\mathrm{SWAP}-\frac Im\right).
\]
The Haar twirl lies in the span of $I$ and $\mathrm{SWAP}$; its trace
is zero and its trace against $\mathrm{SWAP}$ is $\delta_{ib}$, which
fixes these coefficients. Pairing against $F_j\otimes F_a$ yields
\[
 \frac1N\sum_\ell O(U_\ell)_{ji}O(U_\ell)_{ab}
 =\frac{\delta_{ja}\delta_{ib}}p .
\]
A unitary $2$-design reproduces this polynomial moment. Apply this
identity to $\Delta T=T(\Phi)-T(\Psi)$ in \eqref{eq:favg}:
\[
 \|f(\Phi)-f(\Psi)\|_{\rm rms}^2
 =\frac{\norm{\Delta T}^2}{m^2(m+1)^2(m^2-1)}.
\]
Equation~\eqref{eq:choi-norm} proves the claim.
\end{proof}

The frame mechanism is the established unitary-design tomography of
\cite{roth}, with conventions displayed to prevent a norm or dimension
factor from being imported incorrectly.

\begin{corollary}[Positive completion from complete design fidelities]
\label{cor:agf}
If $\Phi$ has Kraus rank at most $r<m$, then its complete exact-design
fidelities identify it among all CPTP maps. For every such competitor,
\[
 \norm{J_\Phi-J_\Psi}
 \le C_{m,r}(m+1)\sqrt{m^2-1}
             \|f(\Phi)-f(\Psi)\|_{\rm rms}.
\]
If $f_{\rm obs}=f(\Phi)+e$ and $\|e\|_{\rm rms}\le\epsilon$, every
minimizer of $\|f(\Psi)-f_{\rm obs}\|_{\rm rms}$ over all CPTP maps obeys
\begin{align}
 \norm{J_{\widehat\Phi}-J_\Phi}
 \le{}&2C_{m,r}(m+1)\sqrt{m^2-1}\,\epsilon \notag\\
 &+\sqrt2(C_{m,r}+1)\tau_r(J_\Phi). \label{eq:agf-noise}
\end{align}
The uniform rank threshold is sharp.
\end{corollary}
\begin{proof}
The exact-rank statement follows from
Theorem~\ref{thm:channel} and Lemma~\ref{lem:agf}.
A residual minimizer has
$\|f(\widehat\Phi)-f(\Phi)\|_{\rm rms}\le2\epsilon$;
apply \eqref{eq:tail} and \eqref{eq:frame} for the approximate-rank
statement. Replacers have rank $m$ and $f_\ell=1/m$ for every
reference unitary, establishing the endpoint obstruction.
\end{proof}

There is no linear reconstruction of the translation in these
fidelities. The additional conclusion is a nonlinear uniqueness and
stability consequence of positivity and rank. The feasible set in
Corollary~\ref{cor:agf} is all CPTP maps; imposing a unital constraint
would discard the nonunital examples above.

\section{Verification, scope and open questions}
The analytic proofs establish statements in all declared dimensions.
Separate finite exact-arithmetic controls check the projection,
attainers, coherent directions, complex Hermitian basis transposes,
rectangular channels and spectral tails. Such controls are diagnostics,
not substitutes for the proofs. The derivation and this combined
manuscript have different review scopes; no score from an earlier
manuscript is assigned to this one.

Three limitations are central. First, the exact constants are for
unrestricted density matrices; optimal constants on the fixed-input-
marginal Choi class remain open. Second, complete-design fidelities
may be numerous. We do not show that a small random Clifford subset
has the required uniform one-sided conditioning for arbitrary
low-Kraus-rank nonunital channels. Third, a deterministic perturbation
bound for calibrated fidelities is not a SPAM-robust experimental
protocol. Noise extraction, gauge and calibration must be justified
separately; the guarantees of \cite{kimmel,roth} do not automatically
transfer to a larger feasible class.

These distinctions suggest concrete further problems: determine the
optimal fixed-marginal constants, find sparse fidelity frames with
uniform positive-completion guarantees, and obtain sharper spectral-tail
terms under additional channel structure. The present paper establishes
the local-erasure law and its complete-data consequence only.

\paragraph{AI and responsibility disclosure.}
The named author directed OpenAI Codex assistance in derivation,
writing and exact arithmetic. Requested GPT-6 Astra xhigh audits
examined component arguments with reused context and shared files.
Runtime identity and isolation were not independently attested.
No independent human review, Lean formalization, certified worldwide
originality or predicted assessment score is claimed. A full assessment
of this exact combined manuscript has not yet been obtained.

\begingroup\small
\begin{thebibliography}{9}
\bibitem{carmeli}
C.~Carmeli, T.~Heinosaari, J.~Schultz and A.~Toigo,
Tasks and premises in quantum state determination,
\emph{J. Phys. A: Math. Theor.} \textbf{47}, 075302 (2014).
\href{https://arxiv.org/abs/1308.5502}{arXiv:1308.5502}.

\bibitem{gutoski}
G.~Gutoski and N.~Johnston,
Process tomography for unitary quantum channels,
\emph{J. Math. Phys.} \textbf{55}, 032201 (2014).
\href{https://doi.org/10.1063/1.4867625}{doi:10.1063/1.4867625}.

\bibitem{kimmel}
S.~Kimmel, M.~P.~da Silva, C.~A.~Ryan, B.~R.~Johnson and T.~Ohki,
Robust extraction of tomographic information via randomized benchmarking,
\emph{Phys. Rev. X} \textbf{4}, 011050 (2014).
\href{https://doi.org/10.1103/PhysRevX.4.011050}{doi:10.1103/PhysRevX.4.011050}.

\bibitem{erratum}
S.~Kimmel et al.,
Erratum: Robust extraction of tomographic information via randomized
benchmarking, \emph{Phys. Rev. X} \textbf{6}, 029902 (2016).
\href{https://doi.org/10.1103/PhysRevX.6.029902}{doi:10.1103/PhysRevX.6.029902}.

\bibitem{roth}
I.~Roth, R.~Kueng, S.~Kimmel, Y.-K.~Liu, D.~Gross, J.~Eisert and M.~Kliesch,
Recovering quantum gates from few average gate fidelities,
\emph{Phys. Rev. Lett.} \textbf{121}, 170502 (2018).
\href{https://doi.org/10.1103/PhysRevLett.121.170502}{doi:10.1103/PhysRevLett.121.170502}.
\end{thebibliography}
\endgroup
\end{document}
